% ============================================================
% Chapter 4: Proposed Method: Specification and Design
%   Section: The Fact-Checking Pipeline
%     Subsection: Article-Level Aggregation
% Included from chapters/04-proposed-method/06-fact-checking-pipeline/fact-checking-pipeline.tex
% Source: docs/arquitectura-tecnica.md 2.4.7; fact_checker.py
% (_VERDICT_SEVERITY).
% ============================================================

\subsection{Article-Level Aggregation}
\label{subsec:aggregation}

An article's verdict is neither an average nor the verdict of its most
important claim: it is the worst of its claims' verdicts, in a fixed
order of severity,
\[
  \texttt{TRUE} < \texttt{PARTIALLY\_TRUE} < \texttt{UNVERIFIED} <
  \texttt{MISLEADING} < \texttt{FALSE}.
\]
\texttt{PARTIALLY\_TRUE} sits below \texttt{UNVERIFIED} because a
claim whose detail is off has at least been checked. The article's
confidence is the mean of its claims' confidences. An article is
publishable when it was admitted and its verdict is neither
\texttt{FALSE} nor \texttt{MISLEADING}: \texttt{UNVERIFIED} does not
block publication, because with a small local model it is the common
outcome, and the reader view shows each claim's own verdict.

\paragraph{Known limitation.} Because \texttt{UNVERIFIED} outranks
\texttt{TRUE}, a single claim without enough evidence decides the
article's verdict however many others were confirmed, and the mean
confidence mixes confidences of different verdicts. Separating ``how
much could be checked'' from ``what was found'' is left open
(Section~\ref{sec:limitations}). The reader view works around it for
the reader: a claim whose search or model failed reads ``not checked''
and is counted apart, and an article none of whose claims was judged is
shown as not checked rather than as unverified.
